\section{Volume-to-Capacity Analysis}
\label{sec:vc}

\subsection{Methodology}

Volume-to-capacity (V/C) analysis for port connectors requires a modification of standard HPMS methodology. Standard HPMS V/C uses 24-hour AADT divided by hourly capacity (lane count $\times$ 2,200 passenger car equivalents per hour per lane, per HCM 6th Edition \citep{HCM6}) divided by 24 hours. This approach distributes daily traffic uniformly across 24 hours, which is appropriate for mixed-use arterials but not for port connectors.

Port connector traffic is concentrated in a 12-hour operating window (typically 6:00~am to 6:00~pm, coinciding with port gate hours). Within that window, truck arrivals and departures are further concentrated around mid-morning (9:00--11:00~am, peak outbound from port) and mid-afternoon (1:00--4:00~pm, peak inbound to port for the following day's gate opening). The 95th-percentile peak hour on a port connector can carry 10--14\% of daily traffic in a single hour; for a standard highway, the comparable figure is 8--10\%.

The adjusted V/C for peak-hour port connector analysis is:

\begin{equation}
  \text{V/C}_{\text{peak}} = \frac{\text{AADT} \times K_{\text{port}}}{N_{\text{lanes}} \times 2,200 \times \text{PHF}}
  \label{eq:vc}
\end{equation}

where $K_{\text{port}}$ is the port peak-hour fraction (0.12--0.14 for port connectors vs. 0.09--0.10 for standard highways) and PHF is the peak-hour factor (typically 0.92 for freeway conditions). Truck equivalence is applied using PCE = 2.0 for heavy trucks at freeway speeds, per HCM 6th Edition \citep{HCM6}.

Additional calibration evidence from the Port of Savannah (GPA 2023 gate count data \citep{GPA2023}) confirms K$_{\text{port}}$ = 0.13 during peak hours (8--10~am, 2--4~pm), consistent with the POLA/POLB-derived estimate. The Port of Baltimore drayage study (MDOT 2021) reports K$_{\text{port}}$ = 0.11--0.14 across three measurement periods, corroborating the 0.12--0.14 range as an empirical standard for major container port connectors \citep{MDOT_Drayage2021}.

AADT data for each connector segment is retrieved from HPMS 2021 \citep{HPMS2018}. For segments where HPMS AADT is not directly available (three connectors cross port authority property before reaching HPMS coverage), we estimate AADT from port gate count data reported in BTS Port Performance statistics \citep{BTS_airfreight2023} using the following relationship:

\begin{equation}
  \text{AADT}_{\text{connector}} \approx 2 \times \text{Gate trips/day} \times \text{PCE}_{\text{truck}}
\end{equation}

The factor of 2 accounts for both the inbound (empty or pre-positioned container) and outbound (loaded container) truck trips per gate transaction.

\subsection{V/C Results}

Table~\ref{tab:vc} presents V/C ratios for all ten connectors under both the standard HPMS methodology (24-hour AADT basis) and the adjusted peak-hour port methodology.

\begin{table}[h]
\centering
\caption{V/C Analysis: Top 10 US Port Connectors (HPMS 2021 basis)}
\label{tab:vc}
\begin{tabular}{lcccc}
\toprule
Port Connector & AADT & Lanes & V/C (24-hr) & V/C (Peak, Adjusted) \\
\midrule
I-710 Long Beach (LA/LB) & 140,000 & 4--6 & 0.89 & 1.52 \\
I-278 approach (NY/NJ) & 108,000 & 6 & 0.82 & 1.38 \\
I-16/I-95 (Savannah) & 62,000 & 4 & 0.77 & 1.31 \\
I-10/I-610 (Houston) & 198,000 & 8 & 0.94 & 1.22 \\
SR-509/SR-518 (Seattle/Tacoma) & 48,000 & 2--4 & 0.91 & 1.45 \\
I-95/I-695 (Baltimore) & 88,000 & 6 & 0.67 & 1.12 \\
I-64 Hampton Roads (Norfolk) & 74,000 & 4 & 0.77 & 1.28 \\
I-526 Mark Clark (Charleston) & 52,000 & 4 & 0.64 & 1.09 \\
I-395/MacArthur (Miami) & 36,000 & 6 & 0.28 & 0.93 \\
I-880/Adeline (Oakland) & 44,000 & 4 & 0.55 & 1.04 \\
\midrule
\multicolumn{2}{l}{V/C $> 1.2$ count} & & 1 of 10 & 7 of 10 \\
\bottomrule
\end{tabular}
\end{table}

The adjustment from 24-hour to peak-hour analysis changes the finding dramatically: only one connector (I-710) appears oversaturated under the standard 24-hour methodology; seven appear oversaturated under the peak-hour port methodology. This demonstrates that the invisibility of port connectors in standard analysis is a methodological artifact, not a reflection of their actual congestion state.

\subsection{I-710 Long Beach: The Case Study}

The I-710 Long Beach Freeway is the most important 11 miles of road in the US import logistics network. It serves the Ports of Los Angeles and Long Beach combined---the San Pedro Bay port complex---which handles approximately 40\% of all US container imports by TEU volume \citep{BTS_airfreight2023}. Every container from China, Japan, South Korea, and Vietnam that arrives at these ports exits through a truck drayage move on the I-710.

At peak (10:00~am--12:00~pm), the I-710 carries approximately 14,000 trucks per hour in both directions on its 4--6 lane cross-section. At V/C 1.52, it operates in the oversaturated regime: arrival rates exceed departure rates, and queuing is building rather than dissipating. The average truck speed on the I-710 during peak hours is approximately 18--22 mph (source: Caltrans Performance Measurement System, PeMS, for loop detector stations on I-710) versus a free-flow speed of 60 mph.

The capacity constraint of the I-710 is not a uniform bottleneck: the freeway narrows from 6 lanes to 4 lanes at the I-405 interchange (the ``Orange Crush'' complex), where three freight-heavy routes converge. This interchange is the single highest-traffic point in the US truck network; at peak, it processes approximately 42,000 vehicles per hour on an interchange designed for 28,000. The resulting speed drop---from 35 mph on the I-710 approach to 8--12 mph through the interchange---generates the primary dwell cost for port trucks.

A fully loaded container truck traveling the I-710 at peak takes approximately 45--55 minutes to traverse 11 miles to the I-405 interchange. At free-flow speed (60 mph), the transit should take 11 minutes. The average delay per truck trip is therefore approximately 35--44 minutes, or 0.6--0.7 hours. At \$225/hr truck operating cost \citep{ATRI_costs2024} and 35,000 trucks per day:

\begin{equation}
  C_{\text{I-710,daily}} = 35,000 \times 0.65 \text{ hrs} \times \$225 = \$5.1 \text{M/day} = \$1.86 \text{B/year}
\end{equation}

This \$1.86 billion per year in truck delay on a single 11-mile segment is the direct cost of the I-710 bottleneck. It does not include vessel demurrage costs (containers stuck in the port because trucks cannot exit), terminal dwell fees, or downstream supply chain disruption costs.

\subsection{Comparison to T1 Bottleneck Rankings}

The ATRI Top 100 Freight Bottlenecks list \citep{ATRI_costs2024} ranks congestion nodes by truck-hours of delay per year. The top-ranked node in 2023 is the I-285/I-85 interchange in Atlanta, accumulating approximately 1.9 million truck-delay-hours per year over a complex that extends approximately 5 miles. The I-710 port connector, at 35,000 trucks per day $\times$ 0.65 hours delay $\times$ 365 days, accumulates approximately 8.3 million truck-delay-hours per year---four times the top-ranked ATRI bottleneck. It does not appear in the ATRI rankings at all, because ATRI's data coverage excludes the I-710 south of the I-405.

This is not an ATRI methodology error---it reflects the genuine difficulty of instrumenting port approach roads that cross jurisdictional boundaries. But the consequence is that the single most delay-intensive 11-mile road segment in the US freight network is absent from the national freight investment priority list.
